% Prose for this counterexample.  Every region below is delimited by an
% environment that tools/build.py extracts; machine facts (ids, uid, dates,
% urls, enums) live in case.json.  See counterexamples/AGENTS.md.

\begin{cxtitle}
Problem 36: stable Schur positivity
\end{cxtitle}

\begin{cxcredits}
\posedby{Borcea and Br\"and\'en \citeyearpar{AIMStableProblems}}
\foundby{TODO}{TODO}
\formalizedby{S.~Sra}
\auditedby{S.~Sra}
\contributedby{S.~Sra}
\end{cxcredits}

\begin{cxcontext}
We keep the notation of the previous subsection: real stability means nonvanishing when every variable lies in the open upper half-plane; $a_\lambda$ are the coefficients in the Schur expansion $p=\sum_{\lambda\vdash d}a_\lambda s_\lambda$ of a symmetric homogeneous $p$ of degree $d$; $f^\lambda$ counts standard Young tableaux of shape $\lambda$; and $\Imm_\lambda$ is the immanant of a $d\times d$ matrix attached to the irreducible character $\chi^\lambda$, with $\Imm_{(1^d)}=\det$ and $\Imm_{(d)}=\per$.  Problem 36 asks whether stability forces Schur-positivity outright; Problems 37 and 38 ask for the two endpoint inequalities that would bracket the Schur coefficients between the $\det$ and $\per$ ends of the immanant scale.
\end{cxcontext}

% ---------------- result: aim-problem-36 ----------------

\begin{cxsource}{aim-problem-36}
Borcea--Br\"and\'en, AIM problem list \citeyearpar{AIMStableProblems}, Problem 36
\end{cxsource}

\begin{cxstatement}{aim-problem-36}
Let $p$ be a symmetric, real stable and homogeneous polynomial of degree $d$ in $n$ variables with $d\le n$ and suppose that $p$ has at least one positive (and then all non-negative) coefficients in the monomial basis.  Is $p$ Schur-positive?  That is, when $p(z_1,\ldots,z_n)$ is expanded in the Schur-basis
\[p(z_1,\ldots,z_n)=\sum_{\lambda\vdash d}a_\lambda s_\lambda(z_1,\ldots,z_n),\]
is $a_\lambda\ge0$ for all $\lambda\vdash d$?
\end{cxstatement}

\begin{cxsummary}{aim-problem-36}
Borcea--Br\"and\'en Problem 36: stable Schur positivity \citeyearpar{AIMStableProblems}
\end{cxsummary}

\begin{cxcertificate}{aim-problem-36}
Positive-definite $5\times5$ det-pencil; exact Schur expansion with $[s_{(1^5)}]q=-2972$.
\end{cxcertificate}

% ---------------- result: aim-problem-37 ----------------

\begin{cxsource}{aim-problem-37}
Borcea--Br\"and\'en, AIM problem list \citeyearpar{AIMStableProblems}, Problem 37
\end{cxsource}

\begin{cxstatement}{aim-problem-37}
If $A$ is a $d\times d$ matrix then $\det(A)=\Imm_{(1^d)}(A)$ and $\per(A)=\Imm_{(d)}(A)$, where $(1^d)=(1,\ldots,1)$ and $(d)=(1^d)'=(d,0,\ldots,0)$.  Schur's inequality [J,L] asserts that if $A$ is a $d\times d$ positive semi-definite matrix then $\Imm_\lambda(A)\ge f^\lambda\det(A)$ for all $\lambda\vdash d$, where $f^\lambda$ is the number of standard Young tableaux of shape $\lambda$.  A natural real stable extension of this result would be as follows.

\emph{Problem 37.}  Let $p$ be as in Problem 36.  Is $a_\lambda\ge f^\lambda a_{(d)}$ for all $\lambda\vdash d$?
\end{cxstatement}

\begin{cxsummary}{aim-problem-37}
Borcea--Br\"and\'en Problem 37: stable Schur inequality (lower endpoint) \citeyearpar{AIMStableProblems}
\end{cxsummary}

\begin{cxcertificate}{aim-problem-37}
Same pencil; $a_{(1^5)}=-2972<1728=f^{(1^5)}a_{(5)}$.
\end{cxcertificate}

\begin{cxrefutation}
The answer to both is no already at $d=n=5$.  Let
\begingroup
\scriptsize
\[
J_1=
\begin{pmatrix}
38&19&19&2&36\\
19&38&2&19&36\\
19&2&38&19&36\\
2&19&19&38&36\\
36&36&36&36&72
\end{pmatrix},\quad
J_2=
\begin{pmatrix}
38&19&19&17&21\\
19&38&17&19&21\\
19&17&38&34&36\\
17&19&34&38&36\\
21&21&36&36&42
\end{pmatrix},\quad
J_3=
\begin{pmatrix}
38&34&19&17&36\\
34&38&17&19&36\\
19&17&38&19&21\\
17&19&19&38&21\\
36&36&21&21&42
\end{pmatrix},
\]
\[
J_4=
\begin{pmatrix}
38&4&19&2&21\\
4&8&2&4&6\\
19&2&38&4&21\\
2&4&4&8&6\\
21&6&21&6&42
\end{pmatrix},\qquad
J_5=
\begin{pmatrix}
8&4&4&2&6\\
4&38&2&19&21\\
4&2&8&4&6\\
2&19&4&38&21\\
6&21&6&21&42
\end{pmatrix},
\]
\endgroup
and set
\begin{equation}\label{eq:aim36-q}
q(x_1,\ldots,x_5)=\frac1{648}\det\left(\sum_{i=1}^5x_iJ_i\right).
\end{equation}

\begin{theorem}[\statusfalse: stable Schur positivity, AIM Problem 36]\label{thm:aim36}
The polynomial \eqref{eq:aim36-q} is symmetric, homogeneous of degree five in five variables, real stable, and all $126$ of its ordinary monomial coefficients are strictly positive, yet
\[
[s_{(1^5)}]\,q=-2972<0.
\]
\end{theorem}
\begin{proof}
The leading principal minors of $J_i$ are
\[
\begin{array}{c|rrrrr}
 &\Delta_1&\Delta_2&\Delta_3&\Delta_4&\Delta_5\\ \hline
J_1&38&1083&28728&202176&1119744\\
J_2&38&1083&28728&202176&1119744\\
J_3&38&288&8208&202176&1119744\\
J_4&38&288&8208&46656&1119744\\
J_5&8&288&1728&46656&1119744,
\end{array}
\]
all positive, so Sylvester's criterion gives $J_i\succ0$.  If $z_1,\ldots,z_5$ lie in the open upper half-plane and $(\sum_iz_iJ_i)v=0$ for some $v\ne0$, then
\[
0=\operatorname{Im}\left\langle v,\left(\textstyle\sum_iz_iJ_i\right)v\right\rangle=\sum_i\operatorname{Im}(z_i)\langle v,J_iv\rangle>0,
\]
a contradiction; hence $q$ is real stable, and homogeneity of degree five is immediate.  Exact expansion of the determinant gives
\begin{align*}
q={}&1728m_{(5)}+19440m_{(4,1)}+66555m_{(3,2)}+140910m_{(3,1,1)}\\
&+250440m_{(2,2,1)}+547405m_{(2,1,1,1)}+1182860m_{(1^5)},
\end{align*}
which exhibits both the $S_5$-symmetry and the strict positivity of every monomial coefficient.  Inverting the Kostka matrix over $\mathbb Z$ in the partition order $(5),(4,1),(3,2),(3,1,1),(2,2,1),(2,1,1,1),(1^5)$ yields
\begin{align*}
q={}&1728s_{(5)}+17712s_{(4,1)}+47115s_{(3,2)}+56643s_{(3,1,1)}\\
&+62415s_{(2,2,1)}+56437s_{(2,1,1,1)}-2972s_{(1^5)}.
\end{align*}
The last coefficient alone follows from the single integer identity
\begin{align*}
[s_{(1^5)}]q={}&1728-2(19440)-2(66555)+3(140910)\\
&+3(250440)-4(547405)+1182860=-2972.
\end{align*}
The source package recomputes the Schur expansion a second, independent way, by the bialternant formula, and every step uses rational arithmetic only.
\end{proof}

\begin{remark}[Origin of the pencil]
The matrices are not a numerical accident.  Let $V=\ker B\cong S^{(3,2)}$, where $B$ maps the permutation module on $2$-subsets of $[5]$ to $\mathbb R^5$ by $e_{ab}\mapsto e_a+e_b$, let $P_i$ be the rank-three projection attached to the point stabilizer of $i$, and let $G$ be the Gram matrix of the chosen basis of $V$.  Then $J_i=2G(I+5P_i)$, and the whole one-parameter family $\det(\sum_ix_i(I+tP_i))$ has a closed Schur expansion whose $s_{(1^5)}$ coefficient turns negative for $t$ large; $t=5$ clears denominators.  The standalone write-up \path{counterexamples/aim-problem-36/paper.tex} develops this construction in full.
\end{remark}

\begin{theorem}[\statusfalse: lower endpoint, AIM Problem 37]\label{thm:aim37}
The polynomial \eqref{eq:aim36-q} satisfies the hypotheses of Problem 37 but violates its lower endpoint inequality
\[
a_\lambda\ge f^\lambda a_{(d)}.
\]
\end{theorem}
\begin{proof}
Take $\lambda=(1^5)$, so that $f^{(1^5)}=1$.  The inequality would require $a_{(1^5)}\ge a_{(5)}$, that is, $-2972\ge1728$.
\end{proof}

\begin{remark}[Relation to Problem 38]
The same sign obstruction refutes Problem 38 as well: at $\lambda=(5)$ the upper endpoint would require $1728\le-2972$.  That refutation is not carried here, because Theorem~\ref{thm:pot} refutes Problem 38 by a Schur-\emph{positive} polynomial, which is the stronger statement.
\end{remark}
\end{cxrefutation}
